% TOP_PROBLEM: {"problemId": "problem.majority-not-in-acc0", "canonicalTitle": "MAJORITY is not in ACC0", "releaseRank": 266, "primaryDomain": "theoretical_computer_science", "primaryDomainLabel": "Theoretical computer science", "displayStatus": "open", "releaseStatus": "open"}
% !TeX program = lualatex
% Definition and sources only; this file is not an LLM solution attempt.
\documentclass[11pt]{article}
\usepackage[margin=25mm]{geometry}
\usepackage{fontspec}
\setmainfont{Latin Modern Roman}
\usepackage{amsmath,amssymb,mathtools}
\usepackage{unicode-math}
\setmathfont{Latin Modern Math}
\usepackage{xurl,hyperref}
\hypersetup{colorlinks=true,urlcolor=blue}
\setcounter{secnumdepth}{0}
\setlength{\parindent}{0pt}
\setlength{\parskip}{5pt}
\setlength{\emergencystretch}{3em}
\begin{document}
\section*{\texorpdfstring{MAJORITY is not in ACC0}{MAJORITY is not in ACC0}}\label{266}
\subsection{Definitions and mathematical statement}
Let $\operatorname{MAJ}_n(x)=1$ when $\sum_{i=1}^n x_i\ge n/2$. A nonuniform $\mathrm{ACC}^0$ family has fixed constant depth, polynomial size, and gates AND, OR, NOT, and $\mathrm{MOD}_m$, for one fixed integer $m\ge2$. AND, OR and modular gates have unbounded fan-in; a modular gate outputs one when its input sum is divisible by $m$. The question is
\[
 (\operatorname{MAJ}_n)_{n\ge1}\notin\mathrm{ACC}^0\ ?
\]
Thus every possible fixed depth and modulus must be ruled out together with every polynomial-size family. Circuits can be chosen separately at each length. A lower bound for prime modulus alone, or for a fixed size exponent, is not the whole separation.
\par\smallskip\textit{Formulation and concept references:} \href{https://arxiv.org/pdf/2607.12346}{[S1]}.

\subsection{Short English statement}
Can majority be excluded from every constant-depth polynomial-size circuit family built from Boolean and fixed-modulus counting gates?
\subsection{Sources}
\begin{itemize}
\item[S1]Vaibhav Krishan and Sundar Vishwanathan, Degree Lower Bounds for Torus Polynomials and MAJORITY vs ACC0; v1 July14,2026.
\url{https://arxiv.org/pdf/2607.12346}.
\textit{Formulation source.}
\end{itemize}
\end{document}
